\documentclass[12pt,a4paper]{article}
\usepackage[utf-8]{inputenc}
\usepackage[french]{babel}
\usepackage{amsmath}
\usepackage{amssymb}
\usepackage{geometry}
\geometry{left=2cm,right=2cm,top=2cm,bottom=2cm}

\title{Chapitre 3: Mathématiques}
\author{}
\date{}

\begin{document}

\maketitle

\section*{Contenu}

\section*{Chapitre 3: Maluses e£ ecthageen dite}


Oure plus te clous une dpace eachilen ole dinenrtan un.


T Ecutune olaus une Lane athanaruee- D Rl.

Sout E ume BON Ÿe,,.,e.ÿ te E.

Lammme c'e une Loc an peut ae Caul mecs © de E amauesr Aes cegrclauees X = (:) Aaus 
Le Le ë # Æ = Z Xe Ce.


u Æ an œ un cemerphusane Fi) 2 > 5 = 
Has cœume La Lane 2 Æ srthamenmer, em cauneilt foccbement des caerdleumees : 
A CG] <xle42 2 = Z 4 es = L'éælespe, au emcorE x | ‘ 7 FT ele» EE dam cer canons : (E) Es) T 4 2e = Le 5 = X Ce Lay. Lu , <æ a Z 4e (e 7 Zu) 
" 
ds) z. . lee 11 2 ZE x = xT. X = Z xle;? 
24 EX: 
Caumanre deja Vu, LÀ 'esuture clans une BON ex€ fat a aklenx eÆ aux ces ceerxolmnmees Le prctut Pcalune  odlevrveut le fra ol bcalaure uoueR ser @ qui fac cuteruenir La Grama part lras. 

2) Oxthoganal LT Lacan peanctée . 
o St AE, GR) A Cars réf <p Aya -- ip dame A7 : ; jh an ré q—< : PA . up 
e£ À D matruee cnnstituere des uecleures calrnnes de {,..,f dans une BON E. c.e AÀ- Hat, (Qu, ke) 
Dame Lx Loue E : eee kex (AT) 
Ce veut oliee que re vetif, Re her œerdlaumeer X sat olamx lee CAT) 
Puuve : Re, Te A: (ag) ( ps 
ên. 4€ dame 4 14 n À, = _ Z eg e.' am mate Af- k ) Se LE Vet12,.,pt. ) 
"4 ee WŸ > Vas <p fl? -O. 2e —7 ZTace,s X= Y) ee? Va <g <p 4AT X ©, 
= (JC) 

& 
Fais ces AË K- Z Ag %e ef ne B-47 Lol br <a. + A4 X = 2 bye 2€ 4# ex lame. a" &. x € W + > (A LS jee Æy 
_. æ> AT. () o. sé (jen 
On meute Alan vue OW D ax Eh a. amañite" ed ? â dise a Lx Han prarrettis cles anatrutes. 
Rx acleurs , celte Jeter mt lea gras sans sapgeler.

He» mabraites le opt cle Base.


QE, SR-RS Se: Èë et a ET on Lanes d'en cet olsnrtatian à.


On age le re cle. passage le E a° A mabuce PE HR) Pr mat, Cf 2) 
Se 0e Cp a 4, = Z. Pr © eo On rappelle que A æ æ gets cam dame ai 7) olus > et X’ LT ) olams ÿ hn | 
[x ex" 
s De plu ne ot un enoluneyhinne de E de mmabure  hepuces make A lan € et À dus À abs 
A'2 PA P.

Cue re pese -6- E auec les Een.



T BON 4 Matures orthagane lb.


A9 Resultre qoulimimeuce .


Sact E une BON ce E ek EAP AS p uecleurs cle €. De mamens am pue A: Hat, (he). € he R) 
Se Met Ze mature PP canilrtee oles goal tb Hcalrunes 1? alans 
M: ATA (ET 2 CRU) he <fh <<< 
Se A: SR 8-AT- ( bin Ana Ang ba b!


L 
p) Le . fes Bye de x Ce La 7 Ana ) - 4, las cale‘ Er nu Fe er 1° <4:l4; LE CZ Ac , Z - A,j Eà ? z ) Lan <e&le = À de Gp 0e Ê=S 7 = 1 Ac Aux, 
‘ T EN "-4 A = Ca CPE) D ee, ae 
= R.A mes = = box A4 
ZA < 
Meg be ur A & 
u=A 
carp| 

2) Mabure anChagenmbs et changement de are. 
Den tun, Une mabure cannes PE MUR) et PT PT.


On male O7 Ligup mn le ces amalauies @ ayants.


Thecreme « Gue E nr exact euc bol és le elinsesssinis n.

et È ume base PN ele E. Une mource Le. Lane & 
Preuve. Lx ?- Tate af) 
——_— 
PTP. AUD RIT LIL chip > <ei8> Lee <£ 2. > 
be AR 
L'ée À A Fan 41,2 : ; TT tes PTP- (52) 
O Alu Res Mu kl ot une fanuilh de vecteur. et A: Hat, @-f) $ 0 L $ folle on => A'A- Te 
Lee A foule athagnual e ATA clerc Le a ceeffi ets olagasraut 70. 


LE 
Paaporbran : PE Ola) ac e& seulement 42 Les vecteur celenmes Pass Poe ee P arment ane Laure ertlia ane. de R' amuse dur gedet Heufaire ubeeel, » TZ aff de remanquke que: da 
T . + e=s PPT es Vos Clg 
Las 
A A eUET es Le n 7 L pT D a Vis Relee Ft Pa = | . (ÉS Ÿ6 7% 3) Brapuetes de O(a) Ecemaqule (“a “a “4 )eots) ë/à 4 ‘és 
(rare tan OU) &E un aus gare le GL,, @). 
Buutree : 4 Ou) € EUR) : MPE Où) P'P=E hu Pat Cnereraclle et Pt PT 2 Ou) ot tale quan mulbphantant 0 (Q0)EOk)" CO tou) = VTUTUV = VV: I. deme WE OR).

LE, 
u 
3, Ge PE Où) aus P'eom) (PPT) Pass @-P'-PT 
ae -(PT ET PPT PPT. 
Le mabures onthagane les Pacut- done facts a camemex ÊLPT CL came P'orthageuute | on ut qu une matrice 
ot arblaqanale AA ex ueckeurs Lier franc anne RO 
Le MR" 
Prapranc ten: Le VE OU), der VU = +4 (pas une ere, (1) Æ S0h)-$ PEOGU) à cet PAS BE vu saus greugee . 
Col vent au fair que A= et, ) soec (00) -dleto"jolet0 dev)" 

\end{document}
